\scanpage{200}

% B01 / text / top, restart
\deleted{Posons}\inserted{Écrivons} maintenant indiciellement l'application (4) :

% B02 / display / top, (10)
\[
 \mathrm{Et}_{A_0}(s)=E_{\varphi_{A_0}(s)}={}^{s}_{\unsure{0|\vphantom{1}}}R_*(\varphi(s)) \tag*{(10)}
\]

% B03 / text / upper, before (11)
où

% B04 / display / upper, (11)
\[
 \varphi_{\unsure{0A_0|0|A_0}}:S\xrightarrow{\unsure{\sim|\mathrm{su}}}S^* \tag*{(11)}
\]

% B05 / text / upper, bracketing and restart
est l'application \deleted{\unsure{b}} surjective [en fait bijective
\deleted{par}\inserted{vu} l'hypothèse \deleted{faite}\inserted{envisagée plus haut}) \deleted{de} induite par (4), \deleted{\unsure{on}}
$[\![$\deleted{justement définie}
\inserted{Ainsi la formule (9)} s'écrit aussi

% B06 / display / upper middle, (12)
\[
 \forall s\in S,\quad \underset{\substack{\cap\\S^*}}{R_*(s)}\xleftarrow{\sim}E_{\varphi_0(s)}\ .\quad ]\!] \tag*{(12)}
\]

% B07 / text / middle, injectivity
Ainsi, la structure \inserted{de graphe strict} $(S,A)$, $A\subset\unsure{P|\mathcal P|\mathfrak P}_2(S)$, $[\![$satisfaisant
la condition d'injectivité

% B08 / display / middle, (13), left brace and marginal NB
\[
 \left\{\begin{aligned}
 &s\longmapsto\mathrm{Et}_{A_0}(s)\quad\text{de } S\text{ dans }P(S)\text{ injective}\\
 &\text{i.e.}\quad \forall\,s\!,t\in S,\ s\ne t,\quad\exists\,u\in S\text{ tel que l'on ait}\\
 &\qquad(\{s,u\}\in A_0,\{t,u\}\notin A_0)\quad\text{ou l'inverse}\\
 &\qquad(\{t,u\}\in A_0,\{s,u\}\notin A_0)\\
 &\text{i.e.}\quad \forall\,a\in P_2(S),\ \exists\,u\in S,\ s,t\in a,\text{ tels que}\\
 &\qquad\{s,u\}\in A_0,\{t,u\}\notin \unsure{a_0|A_0}\quad ]\!]
 \end{aligned}\right. \tag*{(13)}
\]
% slanted
\margin{\uline{NB} Même sans condition d'injectivité (13), la structure de graphe $(S,A_0)$ définit une structure (14), à cela près que $\varphi_0$ n'est plus néc. bijective, mais seulement surjective. Mais pour la suite cela se trouve être rédhibitoire -- brisant la symétrie entre $S$ et $S'$, en nous empêchant de passer à $\varphi_0^{-1}$ \ldots C'est dès qu'on veut dualiser qu'on est obligé de faire l'hypothèse \unsure{séparants|séparante} \ldots}

% B09 / text / middle, before (14)
définit les structures suivantes

% B10 / display / lower middle, (14), left brace
\[
 \left\{\begin{aligned}
 &\text{a) une partie } S^*\subset P(S),\quad\text{donc un \deleted{\unsure{gr}} ens. }S^*\\
 &\quad\text{et une relation}\\
 &\qquad R_*\subset S\times S^*\\
 &\text{b) une application }[\![\text{bijection}]\!]\\
 &\qquad\varphi_0:S\xrightarrow{\unsure{\sim|\mathrm{su}}}S^*,
 \end{aligned}\right. \tag*{(14)}
\]

% B11 / text / lower, reconstruction
permettant de reconstruire \inserted{la relation \deleted{\unsure{$R$}} $R_0\subset S\times S$, et donc la}
donnée structure de graphe $(S,A_0)$, par la formule

% B12 / display / lower, (15), with gloss underneath
\[
 \forall\,s,t\in S,\quad
 \underset{\substack{[\text{i.e. }s\in\underbrace{\mathrm{Et}_{A_0}(t)}_{={}^s R_0(\varphi_0(t))}]}}{(s,t)\in R_0}
 \Longleftrightarrow \deleted{\unsure{t|s}}(s,\varphi_0(t))\in R_* \tag*{(15)}
\]

% B13 / text / lower, before (15 prime)
i.e.

% B14 / display / lower, (15 prime)
\[
 R_0=(\mathrm{id}_S\times\varphi_0)^{-1}(R_*), \tag*{$(15')$}
\]

% B15 / text / foot, before (16)
d'où

% B16 / display / foot, (16)
\[
 A_0=\left\{a\in P_2(S)\mid\exists\,\unsure{\{s,t\}|s,t|(s,t)}\in\deleted{\unsure{R_0}}S,\ \text{avec }\left\{\begin{aligned}a&=\{s,t\}\\(s,\varphi_0(t))&\in R_*\end{aligned}\right.\right\}. \tag*{(16)}
\]

% archive sheet number: 199
